\chapter{Requisiti e principi di progetto}
\label{cap:03-requisiti-principi}

Il \cref{cap:01-introduzione} ha enunciato la tesi del lavoro: un insieme
selezionato di rinunce produce un sistema pi\`u utile come strumento di misura
di quanto lo sarebbe la sua versione completa. Questo capitolo rende operativa
quella tesi. Per ciascun vincolo dichiara che cosa esclude, che cosa
\emph{abilita}, e in quale punto del documento se ne trova la conseguenza
sperimentale.

\section{I vincoli, e che cosa comprano}

\subsection{Processo di controllo singolo}

\nf{} gira in un solo processo. Non esiste modalit\`a distribuita, non esiste
elezione di leader, non esiste replicazione dello stato di controllo. Il codice
non incapsula questa assunzione dietro un'astrazione: le strutture dati di
coordinamento sono mappe concorrenti in memoria, non registri replicati.

\paragraph{Che cosa esclude.} L'alta disponibilit\`a. Un riavvio del pod
interrompe tutte le invocazioni in volo e azzera il registro delle funzioni. La
piattaforma non \`e adatta a un servizio la cui indisponibilit\`a abbia costo.

\paragraph{Che cosa abilita.} Tre cose distinte, tutte necessarie nei capitoli
sperimentali.

Primo, l'\textbf{osservabilit\`a totale dello stato di controllo}. Con un solo
processo, la profondit\`a di una coda, il numero di richieste in volo per
funzione e il limite di concorrenza corrente sono variabili di un unico spazio
di indirizzamento, leggibili in modo consistente nello stesso istante. In un
control plane distribuito la stessa lettura richiederebbe un'aggregazione, che
introduce ritardo e incertezza proprio nelle grandezze che il
\cref{cap:25-valutazione-concorrenza} deve misurare con risoluzione di
millisecondi.

Secondo, l'\textbf{assenza di dinamiche di coordinamento} nel segnale misurato.
Un controllore di concorrenza replicato deve riconciliare decisioni prese da
istanze diverse su osservazioni parzialmente sovrapposte; le oscillazioni che ne
derivano sono un fenomeno interessante ma \emph{diverso} da quello studiato qui,
e sarebbero indistinguibili da esso.

Terzo, la \textbf{brevit\`a del percorso critico}, quantificata nella
\cref{tab:piattaforme}: una singola invocazione sincrona attraversa un solo
processo di controllo prima di raggiungere il codice utente.

\subsection{Stato in memoria}

Il registro delle funzioni, lo storico delle esecuzioni, le chiavi di
idempotenza e le code sono tutte strutture in memoria. Nulla \`e persistito su
disco o su un sistema di stato esterno.

\paragraph{Che cosa esclude.} La sopravvivenza a un riavvio, e con essa
qualunque garanzia di durabilit\`a sulle invocazioni asincrone accettate ma non
ancora eseguite.

\paragraph{Che cosa abilita.} L'eliminazione dal percorso critico di ogni
operazione di I/O che non sia la chiamata alla funzione stessa. Nelle
piattaforme che consultano un database per risolvere la definizione di una
funzione, quella lettura \`e sul percorso di ogni invocazione; in \nf{} \`e
l'accesso a una \code{ConcurrentHashMap}. La conseguenza \`e che la latenza
misurata dal control plane, al netto dell'esecuzione, \`e dominata da
elaborazione in memoria e non da una coda di I/O il cui comportamento sotto
carico avrebbe dinamiche proprie.

Va notato che questa scelta impone il proprio prezzo anche a regime, non solo al
riavvio: uno store in memoria che non dimentica cresce senza limite. Il
\code{ExecutionStore} (\cref{cap:06-nucleo}) implementa quindi tre orizzonti di
ritenzione distinti --- TTL di conservazione, TTL di alleggerimento del record e
durata massima assoluta --- perch\'e un'esecuzione che non arriva mai a uno stato
terminale non deve poter restare in memoria per sempre.

\subsection{Nessuna autenticazione n\'e autorizzazione}

L'API non autentica i chiamanti e non applica alcun controllo di accesso.

\paragraph{Che cosa esclude.} L'esposizione su una rete non fidata. \nf{} \`e
pensato per girare dentro un perimetro gi\`a protetto --- un cluster di
laboratorio, una VM di test, una rete privata.

\paragraph{Che cosa abilita.} Meno di quanto si potrebbe pensare sul piano delle
prestazioni: il costo di una verifica di token \`e reale ma modesto. Il guadagno
\`e principalmente sulla superficie di codice e sul numero di modi in cui un
esperimento pu\`o fallire per ragioni non attinenti all'esperimento. Un carico di
prova che riceve \code{401} per un token scaduto \`e tempo perso in un contesto
in cui l'unica cosa che si vuole osservare \`e la distribuzione della latenza.

Questa \`e la rinuncia meno difendibile delle tre, e il \cref{cap:27-discussione}
la riprende: \`e la sola che impedisce alla piattaforma di essere usata al di
fuori di un laboratorio, ed \`e anche la sola per cui l'aggiunta non
comprometterebbe nessuno degli obiettivi sperimentali.

\subsection{Prestazioni e latenza prima delle funzionalit\`a}

Dove una funzionalit\`a aggiuntiva comporterebbe lavoro sul percorso di
invocazione, la funzionalit\`a viene rifiutata o spostata fuori dal percorso.

Questo principio ha conseguenze visibili nel codice. Il \code{ExecutionStore}
espone due varianti di lettura, una che restituisce \code{Optional} e una che
restituisce \code{null}, con il commento che la seconda \`e destinata al percorso
caldo per evitare l'allocazione dell'\code{Optional}. Il metodo che risolve
l'idempotenza pu\`o attendere brevemente su una rivendicazione contesa, e per
questo \`e esplicitamente confinato su uno scheduler elastico anzich\'e sul
thread dell'event loop Netty. Sono micro-decisioni; ci\`o che conta \`e che
esistano e siano commentate come tali, perch\'e definiscono il criterio con cui
il progetto arbitra i propri compromessi.

\section{Il principio architetturale: nucleo minimale e moduli}

Il vincolo pi\`u caratterizzante non \`e una rinuncia ma una struttura: il
control plane \`e diviso in un \textbf{nucleo} che implementa i quattro obblighi
del \cref{cap:02-stato-dell-arte} e in \textbf{moduli opzionali} che ne
estendono il comportamento.

La regola che governa la divisione \`e la seguente:

\begin{quote}
Il nucleo deve poter essere compilato ed eseguito senza alcun modulo, e in
quella configurazione deve restare una piattaforma FaaS funzionante --- ridotta,
ma non rotta.
\end{quote}

\`E una regola pi\`u forte di quanto sembri. Non dice che i moduli sono
disattivabili con un flag: dice che il binario prodotto senza moduli non li
contiene affatto. La selezione avviene a tempo di compilazione
(\cref{cap:08-sistema-moduli}), e quindi il codice del nucleo non pu\`o
contenere rami condizionali che testano la presenza di un modulo. Deve invece
dipendere da \emph{interfacce} di cui esiste sempre almeno un'implementazione
neutra.

\subsection{I punti di estensione e i loro default neutri}

Il nucleo dichiara quattro interfacce e fornisce per ciascuna
un'implementazione no-op, attiva quando nessun modulo ne registra una:

\begin{table}[htbp]
\centering
\caption{Punti di estensione del nucleo e comportamento senza moduli.}
\label{tab:estensioni}
\small
\begin{tabularx}{\textwidth}{@{}lXX@{}}
\toprule
\textbf{Interfaccia} & \textbf{Fornita da} & \textbf{Comportamento del default} \\
\midrule
\code{InvocationEnqueuer} & \code{async-queue} &
Dichiara \code{enabled() == false}; il percorso asincrono risponde
\code{501 Not Implemented}, il sincrono dispaccia in linea. \\
\code{SyncQueueGateway} & \code{sync-queue} &
Nessuna ammissione: ogni richiesta sincrona procede senza controllo di
backpressure. \\
\code{ScalingMetricsSource} & \code{async-queue} &
Restituisce zero per profondit\`a e richieste in volo, e dichiara di non essere
una sorgente reale. \\
\code{ImageValidator} & provider di deployment &
Accetta ogni immagine senza verificarne l'esistenza. \\
\bottomrule
\end{tabularx}
\end{table}

Il quarto caso merita una nota, perch\'e non segue lo schema degli altri tre. La
validazione delle immagini non \`e un modulo autonomo: ciascun provider di
deployment possiede il proprio validatore
(\code{KubernetesImageValidator}, \code{DockerImageValidator}) e lo attiva
quando viene selezionato come backend. La ragione \`e che verificare l'esistenza
di un'immagine \`e un'operazione specifica del backend --- interrogare un
registry attraverso l'API Kubernetes non \`e la stessa cosa che interrogare un
demone Docker locale --- e un modulo generico avrebbe dovuto reintrodurre la
distinzione al proprio interno.

\subsection{Quando un default neutro non \`e accettabile}

Il principio ha un'eccezione codificata, ed \`e istruttiva. Il modulo
\code{concurrency-control} legge profondit\`a di coda e richieste in volo da
\code{ScalingMetricsSource}, e --- soprattutto --- \emph{riscrive} attraverso la
stessa interfaccia i limiti che calcola. \`E quella scrittura a rendere i limiti
effettivi. Selezionato senza \code{async-queue}, il modulo leggerebbe zeri e
scriverebbe verso un'implementazione che scarta: girerebbe, produrrebbe metriche
plausibili, e non governerebbe nulla.

Il modulo quindi \textbf{rifiuta di avviarsi} e nomina il modulo mancante. La
distinzione che il progetto traccia \`e fra degradazione \emph{osservabile} e
degradazione \emph{silenziosa}: la prima \`e accettabile, la seconda no. Lo
stesso criterio spiega perch\'e il modulo \code{autoscaler}, che legge la stessa
sorgente ma la cui metrica principale (\code{rps}) proviene da un contatore
indipendente, si limita invece a registrare un avviso quando una delle sue
metriche derivate diventa cieca.

\section{Requisiti funzionali}

\begin{description}[style=nextline, leftmargin=0pt]
  \item[\textbf{RF1 --- Registrazione e ciclo di vita.}] La piattaforma deve
  consentire di registrare, aggiornare, elencare, ispezionare e cancellare una
  funzione tramite API REST, con validazione della specifica al momento della
  registrazione (\cref{cap:05-contratti-dati}).

  \item[\textbf{RF2 --- Invocazione sincrona e asincrona.}] Deve offrire
  un'invocazione sincrona che restituisce l'esito nella stessa risposta HTTP e
  un'invocazione asincrona che restituisce un identificatore di esecuzione
  consultabile in seguito (\cref{cap:06-nucleo}).

  \item[\textbf{RF3 --- Idempotenza.}] Due invocazioni che portano la stessa
  chiave di idempotenza sulla stessa funzione devono corrispondere a una sola
  esecuzione, e la seconda deve poter attendere o riprodurre l'esito della prima
  (\cref{cap:06-nucleo}).

  \item[\textbf{RF4 --- Modalit\`a di esecuzione multiple.}] Deve supportare
  l'esecuzione in-process (per il collaudo), l'inoltro a un endpoint esterno gi\`a
  esistente, e la gestione del ciclo di vita di un deployment attraverso un
  backend sostituibile (\cref{cap:07-modalita-esecuzione}).

  \item[\textbf{RF5 --- Riprova e timeout.}] Deve applicare un timeout per
  invocazione e un numero massimo di tentativi, entrambi configurabili per
  funzione con default di piattaforma.

  \item[\textbf{RF6 --- Osservabilit\`a.}] Deve esporre metriche in formato
  Prometheus su una porta separata da quella dell'API, con granularit\`a per
  funzione su latenza, esito, cold start, profondit\`a di coda e concorrenza
  effettiva (\cref{cap:21-osservabilita}).
\end{description}

\section{Requisiti non funzionali}

\begin{description}[style=nextline, leftmargin=0pt]
  \item[\textbf{RNF1 --- Componibilit\`a verificata.}] Ogni combinazione di moduli
  dichiarata valida deve essere costruibile e la suite di test deve poter girare
  su di essa. La configurazione priva di moduli deve avere una copertura di test
  propria (\cref{cap:16-moduli-diagnostici}).

  \item[\textbf{RNF2 --- Isolamento fra moduli.}] Un modulo non deve dipendere da
  un altro modulo se non attraverso interfacce dichiarate nel nucleo o nel
  modulo condiviso dei contratti. La regola \`e verificata automaticamente
  (\cref{cap:26-qualita-verifica}).

  \item[\textbf{RNF3 --- Non bloccante sul percorso di invocazione.}] Il percorso
  sincrono non deve bloccare thread dell'event loop. Le operazioni
  potenzialmente bloccanti devono essere esplicitamente confinate su scheduler
  dedicati (\cref{cap:04-architettura-control-plane}).

  \item[\textbf{RNF4 --- Compilabilit\`a in immagine nativa.}] L'intero control
  plane, con tutti i moduli, deve essere compilabile con GraalVM Native Image
  senza rinunciare a funzionalit\`a, il che vincola l'uso della riflessione e
  richiede la dichiarazione esplicita degli \emph{hint} di runtime
  (\cref{cap:20-build-packaging}).

  \item[\textbf{RNF5 --- Riproducibilit\`a delle misure.}] Ogni numero
  prestazionale pubblicato deve essere associato a un profilo hardware
  identificato e a un record versionato, e confrontabile con quello della
  release precedente sullo stesso profilo (\cref{cap:22-metodo-sperimentale}).
\end{description}

\section{I default di piattaforma}

I valori di default definiscono il comportamento della piattaforma in assenza di
configurazione esplicita, e vale la pena esaminarli perch\'e ne rivelano le
priorit\`a. Sono dichiarati in \code{application.yml}:

\begin{center}
\small
\begin{adjustbox}{max width=\textwidth}
\begin{tabular}{@{}llp{7.2cm}@{}}
\toprule
\textbf{Chiave} & \textbf{Default} & \textbf{Interpretazione} \\
\midrule
\code{defaults.timeoutMs} & 30\,000 & Generoso: pensato per non far fallire un
esperimento a causa del timeout. \\
\code{defaults.concurrency} & 4 & Conservativo: quattro invocazioni contemporanee
per funzione. \\
\code{defaults.queueSize} & 100 & Coda per funzione sul percorso asincrono. \\
\code{defaults.maxRetries} & 3 & \\
\code{rate.maxPerSecond} & 1\,000\,000 & Di fatto disattivato: il limitatore
esiste ma il default non vincola. \\
\code{sync-queue.max-depth} & 200 & \\
\code{sync-queue.max-estimated-wait} & 2\,s & Soglia di rifiuto per attesa
stimata. \\
\code{scaling.poll-interval-ms} & 5\,000 & Periodo di campionamento
dell'autoscaler. \\
\code{metrics.profile} & \code{basic} & Il profilo esteso \`e opt-in. \\
\code{admin.runtime-config.enabled} & \code{false} & L'API di riconfigurazione a
caldo \`e disattivata per default. \\
\bottomrule
\end{tabular}
\end{adjustbox}
\end{center}

Due default meritano commento. Il limite di frequenza a un milione di richieste
al secondo \`e un limitatore \emph{presente ma neutro}: il meccanismo esiste ed
\`e sul percorso, cos\`i che il suo costo sia incluso in ogni misura, ma non
respinge nulla finch\'e un operatore non lo configura. \`E l'opposto della scelta
usuale, e discende dal principio che un esperimento non deve fallire per una
protezione che non stava studiando. Simmetricamente, l'API di riconfigurazione a
caldo \`e disattivata per default perch\'e espone operazioni di scrittura su un
sistema che non autentica: attivarla \`e una decisione che l'operatore deve
prendere consapevolmente.
